\section{Integration with the BX3 Framework}
\label{sec:rs-integration}

The Recursive Spawning Protocol is a structural instantiation of the BX3 three-layer accountability architecture, not a module composed alongside it as an incidental component. Every RSP mechanism maps onto exactly one BX3 layer, and this mapping is minimal and structurally necessary: no RSP guarantee can be achieved with fewer than all three layers, and no layer can subsume another's role without collapsing at least one correctness property. This section specifies the three mappings and demonstrates why each is exclusive to its layer.

% -------------------------------------------------------
\subsection{Purpose Layer: Accountability Anchor}
\label{sec:rs-purpose-anchor}

The Purpose Layer occupies two structurally singular roles that no other BX3 layer can perform. Both are constitutive for RSP's correctness properties.

\medskip
\noindent\textbf{Exclusive origin of spawn authorization.} At root provisioning, the Purpose Layer defines and records the spawn authorization policy $P_{\mathrm{spawn}}$ in the Fact Layer authorization ledger. $P_{\mathrm{spawn}}$ encodes two mandatory pre-conditions for any child spawn: (i) the child operating scope must be a strict descendant refinement of the parent scope ($R(n_{i+1}) \subset R(n_i)$), and (ii) the spawn must be operationally necessary --- the unmet condition must be unsatisfiable within the parent's existing scope. These conditions originate exclusively at the Purpose Layer. The Bounds Engine evaluates proposed spawns against $P_{\mathrm{spawn}}$; the Fact Layer enforces it; neither originates it. No machine actor may issue a spawn warrant in the absence of a matching Purpose Layer authorization record. The separation is architectural, not organizational: it holds regardless of privilege level, operational urgency, or system state.

\medskip
\noindent\textbf{Exclusive terminus of escalation.} When the chain reaches $D_{\mathrm{ESCALATE}}$, the protocol compulsory-escalates to the human anchor $h(n)$. Critically, $h(n)$ is non-collapsing: for all nodes $n_i$ in chain $C$, $h(n_i) = h(n_0)$. The human anchor does not migrate down the chain as depth increases. The human issues one of three directives --- \texttt{ACK}, \texttt{RESOLVE}, or \texttt{REJECT} --- and no machine actor may issue any of these. The chain terminates in human judgment, never in autonomous resolution. This is not a policy convention. It is the structural expression of the upstream BX3 accountability guarantee: systems built on the BX3 Framework fail upward into human consciousness, never downward into autonomous chaos.

\medskip
The structural necessity of the Purpose Layer's exclusivity is immediate. If any machine actor could absorb an exception or issue a terminal directive, the accountability chain would terminate at that actor rather than at a human. The system would compound failure rather than surface it. The Purpose Layer's exclusivity here is load-bearing: removing it causes the upstream accountability guarantee to collapse into a machine-actor loop.

% -------------------------------------------------------
\subsection{Bounds Engine: Depth Enforcement}
\label{sec:rs-bounds-depth}

The Bounds Engine provides constrained execution evaluation: it evaluates proposed actions against the operating range defined by the Purpose Layer and enforces the depth constraint that prevents unbounded chain growth. Critically, the Bounds Engine cannot exceed these constraints under any operational circumstance --- there is no privilege level, operational exigency, or system state under which it may override the depth bound.

\medskip
\noindent\textbf{Depth evaluation at spawn.} At each spawn event, the Bounds Engine evaluates current chain depth $|C|$ against the Purpose-Layer-defined maximum $D_{\mathrm{max}}$, both anchored in the Fact Layer authorization ledger. If $|C| \geq D_{\mathrm{max}}$, the Bounds Engine does not execute the spawn. It transitions to \texttt{ESCALATING} state and initiates the Bailout Protocol, propagating the exception upward to $h(n)$. This enforcement is deterministic: the Bounds Engine has no discretion to exceed $D_{\mathrm{max}}$ regardless of urgency, operational exigency, or the importance of the unmet condition. The depth constraint is a hard architectural boundary, not a heuristic or a default overridable under pressure.

\medskip
\noindent\textbf{Scope refinement evaluation.} The Bounds Engine also evaluates the scope subset condition $R(n_{i+1}) \subset R(n_i)$ at each spawn event, determining whether a spawn request is justified by operational necessity within the Purpose Layer's authorization scope. This evaluation is a precondition for Fact Layer pre-commit hook invocation; it is not sufficient alone, because the scope refinement condition is independently enforced by the Fact Layer. The separation reflects the depth-limit insufficiency theorem: depth limits alone cannot prevent scope creep drift, so both depth enforcement and scope refinement enforcement are required as structurally distinct operations.

\medskip
\noindent\textbf{Structural reinforcement by the Fact Layer.} The Bounds Engine's depth evaluation is structurally reinforced by the Fact Layer pre-commit hook. Before any spawn operation is recorded, the Fact Layer independently verifies two conditions: (a) $R(n_{i+1}) \subseteq R(n_i)$ (scope subset), and (b) $\mathrm{depth}(n_{i+1}) \leq D_{\mathrm{max}}$ (depth bound). The Fact Layer rejects any spawn violating either condition. The Bounds Engine cannot execute a rejected spawn; it may only transition to \texttt{ESCALATING}. The pre-commit hook makes the depth bound and scope refinement constraint structural invariants --- not policy conventions overridable in exigent circumstances.

The structural necessity of the Bounds Engine follows from its role as the layer responsible for evaluating whether a spawn is operationally necessary within the authorized scope and for managing the transition to \texttt{ESCALATING} state when the depth bound is reached. The Fact Layer enforces what the Bounds Engine evaluates; neither layer alone is sufficient.

% -------------------------------------------------------
\subsection{Fact Layer: Forensic Ledger}
\label{sec:rs-fact-ledger}

The Fact Layer provides deterministic enforcement and forensic auditability. It maintains the append-only ledger that makes the RSP chain auditable at any future time and enforces the write protocol constraints that make RSP structurally tamper-evident. The Fact Layer is the structural foundation on which both the Purpose Layer's authorization policy and the Bounds Engine's depth constraints become runtime-enforceable, not merely documented.

\medskip
\noindent\textbf{Append-only forensic ledger.} The Fact Layer records every spawn event, every state transition, every Purpose Layer directive received, and every unresolved exception condition. Each entry is timestamped, attributed to the node that generated it, and hash-chained to the preceding entry. Hash-chaining ensures that no entry can be inserted, deleted, reordered, or altered after recording: any tampering breaks hash continuity and is detectable by any observer with ledger read access. The ledger is the single source of truth for all authorized chain state --- the runtime artifact that distinguishes RSP chains from conventional delegation chains, which exist only as forensic reconstructions attempted after the fact.

\medskip
\noindent\textbf{Immutable parent pointer.} Once $\alpha(n)$ is established at node provisioning, no actor --- including the Purpose Layer after provisioning --- may modify it. The Fact Layer write protocol rejects any attempt to overwrite $\alpha(n)$. Chain topology changes occur only through authorized spawn operations recorded in the ledger; there is no mechanism for retroactive pointer manipulation. The parent pointer is not merely access-controlled --- it is structurally immutable by protocol definition. This eliminates the re-parenting drift pathway by making retroactive chain rewriting structurally impossible regardless of the requestor's identity or privilege level.

\medskip
\noindent\textbf{Pre-commit hook.} The pre-commit hook is the enforcement mechanism that makes all other RSP constraints structural. At every spawn event, it verifies conditions (a) and (b) before the spawn is recorded. Rejected spawns are logged; the Bounds Engine does not execute them; the protocol transitions to \texttt{ESCALATING} if the unmet condition persists. The pre-commit hook operates identically on all machine actors; there is no privileged actor capable of bypassing it. Even administrative compromise of a machine actor cannot produce a bypass of the pre-commit hook, because the hook is enforced at the Fact Layer write protocol level, not at the actor's discretion.

The structural necessity of the Fact Layer follows from the immutable-parent-pointer requirement: without the Fact Layer write protocol, immutability is a promise rather than a constraint. A mutable parent pointer enables all drift failure modes. Without the pre-commit hook, the scope refinement constraint and depth bound are policy documentation --- enforceable only until an actor with sufficient privilege decides they are inconvenient. The Fact Layer is the structural enforcement layer that makes the Purpose Layer's policy and the Bounds Engine's constraints operative.

% -------------------------------------------------------
\subsection{Structural Minimality}
\label{sec:rs-minimality}

The three-layer mapping is minimal and sufficient. Removing any layer causes the corresponding guarantee to fail, as shown in Table~\ref{tab:rs-layer-minimality}.

\begin{table}[h]
\begin{center}
\begin{tabular}{p{3.4cm} p{4.6cm}}
\toprule
\textbf{Removed layer} & \textbf{Guarantee that fails} \\
\midrule
Purpose Layer & Exclusive human terminus eliminated; chain can terminate in machine actors; spawn policy becomes self-authorized by machine actors \\
Bounds Engine & Constrained execution evaluation eliminated; no layer to evaluate spawn necessity or manage escalation transition \\
Fact Layer & Immutable parent pointer eliminated; retroactive chain manipulation possible; scope refinement and depth bound become policy conventions \\
\bottomrule
\end{tabular}
\caption{Layer removal consequences for RSP correctness properties.}
\label{tab:rs-layer-minimality}
\end{center}
\end{table}

The three-layer architecture is the minimal structure that makes RSP's correctness properties unconditionally guaranteed. No subset of layers achieves the same guarantees.

% -------------------------------------------------------
% SECTION 7
% -------------------------------------------------------
\section{Correctness Properties}
\label{sec:rs-correctness}

This section establishes three correctness properties for the Recursive Spawning Protocol: drift prevention, chain integrity, and termination. Each property is stated formally, proved sketchedly, and connected to the BX3 Framework layer that enforces it. Together these properties constitute the RSP's overall guarantee: a spawn chain is an accountable, bounded, and auditable refinement process that terminates in human judgment.

% -------------------------------------------------------
\subsection{Drift Prevention}
\label{sec:rs-drift-prevention}

\medskip
\noindent\textbf{Theorem (Drift Prevention).} Let $C = \langle n_0, n_1, \ldots, n_k \rangle$ be a Recursive Spawning chain rooted at $n_0$ with human anchor $h(n_0) = h(n_k)$. Then for all $i$ ($0 \leq i \leq k$), the operating scope $R(n_i)$ is a descendant refinement of the intent of $h(n_0)$. Equivalently: no node in $C$ can exhibit behavior outside the authorized intent of its human anchor.

\medskip
\noindent\textbf{Proof sketch.} The proof proceeds by structural induction on chain depth, where each inductive step is enforced at runtime by the Fact Layer pre-commit hook.

\medskip
\textit{Base case ($i=0$).} $R(n_0)$ is defined by $h(n_0)$ at provisioning. By the Purpose Layer's definition of authorized operating scope, $R(n_0)$ is a descendant refinement of $h(n_0)$'s intent by construction. No drift exists at the root.

\medskip
\textit{Inductive step.} Assume $R(n_i)$ is a descendant refinement of $h(n_0)$'s intent for some $i$ with $0 \leq i < k$. For $n_i$ to spawn $n_{i+1}$, the Fact Layer pre-commit hook verifies two conditions before the spawn is recorded:

\begin{enumerate}
    \item[(a)] $R(n_{i+1}) \subseteq R(n_i)$ --- the scope refinement constraint. The Purpose Layer's spawn authorization policy $P_{\mathrm{spawn}}$ requires strict scope refinement; there is no authorized provision for lateral or upward scope expansion.
    \item[(b)] $\mathrm{depth}(n_{i+1}) \leq D_{\mathrm{max}}$ --- the depth constraint, also enforced by the pre-commit hook.
\end{enumerate}

The Fact Layer rejects any spawn operation violating (a) or (b). No actor, including the Bounds Engine, can execute a rejected spawn. Therefore $R(n_{i+1}) \subseteq R(n_i)$ holds for the child by structural enforcement. By the inductive hypothesis, $R(n_i)$ is a descendant refinement of $h(n_0)$'s intent. Since $R(n_{i+1}) \subseteq R(n_i)$, transitivity of subset gives that $R(n_{i+1})$ is also a descendant refinement of $h(n_0)$'s intent. Hence $n_{i+1}$ does not drift. By induction, no node in $C$ drifts. $\square$

\medskip
The proof sketch rests on a conjunction of three architectural facts that distinguish RSP drift prevention from conventional policy-based arguments:

\begin{enumerate}
    \item[(1)] \textbf{Immutable parent pointer.} $\alpha(n)$ is established once at provisioning and is thereafter immutable under the Fact Layer write protocol. Scope inheritance is traceable to a unique authorized lineage, not a manipulable graph.
    \item[(2)] \textbf{Forensic ledger.} Every scope assignment is recorded in the append-only, hash-chained ledger. The ledger makes the inductive proof empirically verifiable: each step of the refinement chain is a recorded, tamper-evident ledger entry.
    \item[(3)] \textbf{Chain terminates at Purpose Layer.} The exclusive human terminus closes the inductive argument at a finite bound. There is no machine-actor alternative terminus.
\end{enumerate}

If any of these three conditions were absent, drift prevention would not be structurally guaranteed. With all three, drift is architecturally impossible.

% -------------------------------------------------------
\subsection{Chain Integrity}
\label{sec:rs-chain-integrity}

\medskip
\noindent\textbf{Theorem (Chain Integrity).} For any Recursive Spawning chain $C = \langle n_0, n_1, \ldots, n_k \rangle$, the parent pointer sequence $\alpha(n_i) = n_{i-1}$ holds for all $1 \leq i \leq k$, and no node configuration in $C$ has been modified after provisioning.

\medskip
\noindent\textbf{Proof sketch.} The proof has two components: topological integrity and node immutability.

\medskip
\textit{Topological integrity.} Each spawn event $n_{i-1} \xrightarrow{\mathrm{spawn}} n_i$ establishes $\alpha(n_i) = n_{i-1}$ as an immutable Fact Layer property at provisioning time. The Fact Layer write protocol rejects any operation attempting to modify $\alpha(n_i)$ after provisioning. Since the chain is constructed exclusively through authorized spawn events and no parent pointer is modified after establishment, $\alpha(n_i) = n_{i-1}$ holds for all $i \geq 1$ by construction. The parent pointer is not merely access-controlled --- it is structurally immutable. There is no privileged actor, including the Purpose Layer after provisioning, capable of overwriting $\alpha(n)$. This eliminates the re-parenting drift pathway by making retroactive chain rewriting structurally impossible.

\medskip
\textit{Node immutability.} Each node $n_i$ is provisioned with a fixed configuration $(R(n_i), P_{\mathrm{spawn}}(n_i), h(n_i))$ established by the Purpose Layer and recorded in the Fact Layer authorization ledger. The Fact Layer enforces immutability of these parameters after provisioning. The write protocol rejects any modification request from any actor --- including the Bounds Engine and the Purpose Layer after provisioning. The Purpose Layer may issue a \texttt{RESOLVE} directive that redirects a node's behavior, but this directive is recorded as a new ledger entry (a state transition), not as an alteration of the provisioned configuration.

The forensic ledger's hash-chaining provides evidentiary proof of chain integrity: any attempt to insert a falsified spawn event or reorder existing entries breaks hash continuity and is detectable by any observer with ledger read access. Chain integrity is not merely a property of the current chain state --- it is a property of the entire historical record, enforceable at any future time. $\square$

% -------------------------------------------------------
\subsection{Termination}
\label{sec:rs-termination}

\medskip
\noindent\textbf{Theorem (Termination).} Every Recursive Spawning chain $C = \langle n_0, n_1, \ldots, n_k \rangle$ terminates --- either at a resolution state or at the escalation threshold --- within a finite number of spawn steps bounded by $\max(D_{\mathrm{max}}, D_{\mathrm{ESCALATE}})$.

\medskip
\noindent\textbf{Proof sketch.} The proof establishes two lemmas and combines them.

\medskip
\textit{Lemma 1 (Depth-bounded growth).} The chain grows only through spawn events. By the Bounds Engine depth enforcement and the Fact Layer pre-commit hook, any spawn event producing $|C| \geq D_{\mathrm{max}}$ is rejected. Therefore the chain cannot exceed depth $D_{\mathrm{max}} - 1$ through autonomous spawn. Rejected spawns are logged; the Bounds Engine transitions to \texttt{ESCALATING}. Hence autonomous chain growth is bounded by $D_{\mathrm{max}}$.

\medskip
\textit{Lemma 2 (Escalation-triggered halt).} If the chain reaches $D_{\mathrm{ESCALATE}}$ before resolving the exception condition, the protocol enters \texttt{ESCALATING} state and suspends all autonomous spawn activity. The Bounds Engine enters quiescent hold, awaiting a directive from $h(n_k)$ --- the human anchor. No further spawn events occur until $h(n_k)$ issues \texttt{ACK}, \texttt{RESOLVE}, or \texttt{REJECT}. Since $h(n_k) = h(n_0)$ is a human principal, the directive arrives in finite time. The \texttt{REJECT} directive terminates the chain definitively; \texttt{RESOLVE} redirects behavior and records a new ledger entry; \texttt{ACK} permits continued operation under continued Purpose Layer authorization.

\medskip
\textit{Theorem conclusion.} Combining Lemma 1 and Lemma 2: the chain grows autonomously at most $D_{\mathrm{max}} - 1$ spawn steps. If it reaches $D_{\mathrm{ESCALATE}}$ before resolution, it halts and escalates. If it reaches $D_{\mathrm{max}}$ before resolution, the Fact Layer rejects further spawn and triggers mandatory escalation. In all execution paths, the chain reaches a terminal state within $\max(D_{\mathrm{max}}, D_{\mathrm{ESCALATE}})$ spawn steps. Since both $D_{\mathrm{max}}$ and $D_{\mathrm{ESCALATE}}$ are finite integers defined at provisioning, the chain terminates in finite time. $\square$

The termination guarantee also precludes infinite spawn-escalate cycles. When $D_{\mathrm{max}}$ is reached, the Fact Layer blocks further spawn and triggers mandatory escalation. The human anchor's \texttt{REJECT} directive terminates the chain. There is no mechanism by which the chain can circumvent this bound and re-enter autonomous spawn --- the Fact Layer pre-commit hook is structural, not configurable at runtime by any machine actor.

% -------------------------------------------------------
\subsection{Collective Guarantee}
\label{sec:rs-collective-guarantee}

The three correctness properties are mutually reinforcing and jointly enforced by the BX3 Framework's three-layer architecture. Together they constitute the RSP's overall accountability guarantee in the context of recursive agent population dynamics.

\medskip
\textbf{Drift prevention} ensures that scope refinement is the only authorized direction of chain growth. Each child scope is a strict subset of its parent, traceable to the human anchor's intent through the forensic ledger. Without drift prevention, successive scope expansions could gradually migrate the chain's operating scope away from the human anchor's intent --- an autonomy runway within an otherwise bounded-depth chain. The depth-limit insufficiency theorem establishes that depth bounds alone cannot prevent this.

\medskip
\textbf{Chain integrity} ensures that the chain topology is stable and tamper-evident. The parent pointer sequence is immutable, node configurations are unmodifiable after provisioning, and the forensic ledger's hash-chaining makes any tampering detectable. Without chain integrity, an adversarial actor could manipulate chain parentage or configuration to redirect accountability or obscure the provenance of a spawned node --- a form of accountability laundering.

\medskip
\textbf{Termination} ensures that the chain cannot grow indefinitely and cannot sustain an infinite spawn-escalate cycle. The depth bound is enforced structurally by the Fact Layer pre-commit hook, and the escalation path terminates in human judgment, never in autonomous resolution. Without termination, an unresolvable exception condition could drive unbounded chain growth exceeding any human's supervisory capacity --- defeating the upstream BX3 accountability guarantee by overwhelming the human anchor.

\medskip
Together these properties guarantee that a Recursive Spawning chain is an accountable, bounded, and auditable refinement process: it grows only in authorized directions (drift prevention), preserves the topology established at provisioning (chain integrity), halts within a Purpose Layer-defined depth bound (termination), and maintains a complete forensic record of every decision point (Fact Layer). This is the operational expression of the upstream BX3 accountability guarantee applied to recursive agent population dynamics.
